\section{生物学示例族 III：节律控制与极限环}

有些系统最好的理解方式，并不是它们会收敛到一个不动点。
相反，它们会生成反复出现的振荡活动。
这也正是极限环与中央模式发生器式推理变得有用的地方。

\begin{definition}[极限环]
\textbf{极限环}是一条稳定的周期轨道。
邻近轨迹会被拉向一个重复运行的轨道，而不是拉向一个静态平衡点。
\end{definition}

\begin{discussion}
节律性生物系统，例如运动控制、类似呼吸的节律以及重复的模式化动作，在教学上非常重要，
因为它们阻止了一个常见错误：并非每一种有用机制都是静止的盆地。
有些机制之所以高产，恰恰是因为它们是循环的。
一旦进入该循环，系统就会以很低的瞬时决策成本通过重复相位持续运转。
\end{discussion}

\begin{example}[把运动常规看作进入一个循环]
运动常规通常不是在重复过程中失败，而是在第一个动作之前失败。
这在数学上很有启发性。
进入该机制可能需要跨越一个很大的势垒；
但一旦课程开始，重复子动作就可能在节律上变得容易。
实践上的教训是：干预设计应当区分\emph{进入循环}与\emph{维持循环}。
这两者不是同一个控制问题。
\end{example}